% Reporte de la Actividad 6
% Sincronizacion de procesos e hilos en Java
\documentclass[12pt,oneside]{template/liverpoolthesis}
\input{template/preamble}

% ---------- Datos de esta actividad ----------
\title{Actividad 6: Sincronización de procesos e hilos en Java}
\author{%
  Andrea Hernandez Huerta\\
  ID: 175523\\[1em]
  Heriberto Espino Montelongo\\
  ID: 175199\\[1em]
  \text{Profesor:} Dr. Victor Giovanni Morales Murillo}
\faculty{Universidad de las Américas Puebla (UDLAP)}
\school{Licenciatura en Ciencia de Datos}
\degree{Curso: Programación Concurrente}
\date{Fecha de entrega: 17 de septiembre de 2026}
\subject{Sincronización de procesos e hilos en Java}

\addto\captionsspanish{\renewcommand{\chaptername}{Ejercicio}}
\providecommand{\figureautorefname}{Figura}
\addto\captionsspanish{\renewcommand{\figureautorefname}{Figura}}

% Inserta una fotografia proporcionada por el estudiante
\newcommand{\espaciofoto}[4]{%
  \clearpage
  \begin{figure}[H]
    \centering
    \IfFileExists{figuras/#1}{%
      \includegraphics[width=0.97\textwidth,height=0.70\textheight,keepaspectratio]{figuras/#1}%
    }{%
      \fbox{%
        \begin{minipage}[c][#4][c]{0.88\textwidth}
          \centering
          \textbf{Espacio reservado para la fotografía}\\[1em]
          #2\\[1em]
          Agregar el archivo como \texttt{\detokenize{figuras/#1}}
        \end{minipage}%
      }%
    }
    \caption{#3}
  \end{figure}%
}

\begin{document}

\frontmatter
\maketitle
\tableofcontents
\listoffigures

\mainmatter

\chapter{Sincronización de hilos con \texttt{join()}}

\textbf{Desarrollen un programa en Java que implemente tres tareas utilizando
\texttt{Runnable}: el Hilo 1 debe simular la obtención de datos, el Hilo 2 el
procesamiento de los datos y el Hilo 3 la generación de un reporte. Utilicen
\texttt{start()} y \texttt{join()} para garantizar el orden Hilo 1, Hilo 2 y
Hilo 3. Cada hilo debe mostrar en consola cuándo inicia y cuándo termina su
tarea.}

\emph{Entradas.} Tres objetos que implementan la interfaz \texttt{Runnable},
uno para obtener datos, otro para procesarlos y otro para generar el reporte.
También se utilizan tres objetos \texttt{Thread} y dos datos compartidos entre
las tareas consecutivas.

\emph{Procesamiento.} El proceso principal crea los tres hilos. Inicia el
Hilo 1 con \texttt{start()} y ejecuta \texttt{join()} para esperar su
finalización. Después inicia el Hilo 2 y vuelve a esperar con \texttt{join()}.
Finalmente inicia el Hilo 3 y espera su terminación. De esta manera, la
finalización de una tarea se convierte en la condición para iniciar la
siguiente.

\emph{Salidas.} Cada hilo muestra un mensaje de inicio y otro de terminación.
El Hilo 1 deja disponibles los datos, el Hilo 2 muestra que fueron procesados
y el Hilo 3 genera el reporte. El mensaje final confirma que las tres tareas
terminaron en el orden solicitado.

\emph{Estructuras de control utilizadas.} Se emplean tres clases que
implementan \texttt{Runnable}, tres objetos \texttt{Thread}, llamadas a
\texttt{start()} para iniciar hilos y llamadas a \texttt{join()} para esperar
su terminación. También se utilizan métodos \texttt{run()}, asignaciones de
datos compartidos y un bloque \texttt{try-catch} para atender una posible
interrupción durante la pausa de cada tarea.

\emph{Conceptos.} En Java, \texttt{Runnable} representa la tarea que puede
ejecutarse en un hilo. \texttt{start()} inicia un nuevo hilo de ejecución y
provoca que la JVM invoque su método \texttt{run()}. \texttt{join()} hace que
el hilo que ejecuta esa instrucción espere hasta que termine el hilo indicado.
En este ejemplo, la espera establece una dependencia: procesar depende de que
la obtención de datos haya terminado y generar el reporte depende de que el
procesamiento haya terminado.

\emph{Prueba de escritorio.} La siguiente secuencia muestra el orden de
ejecución establecido por las llamadas a \texttt{start()} y \texttt{join()}.

\begin{center}
\scriptsize
\begin{tabular}{c|p{3.0cm}|p{8.0cm}}
Paso & Hilo o proceso & Acción y estado \\
\hline
1 & Proceso principal & Crea las tres tareas \texttt{Runnable} y los tres objetos \texttt{Thread}. \\
2 & Hilo 1 & \texttt{start()} inicia la obtención de datos; el hilo muestra que inicia. \\
3 & Proceso principal & Ejecuta \texttt{hilo1.join()} y permanece en espera hasta que el Hilo 1 termina. \\
4 & Hilo 1 & Completa la obtención, guarda los datos y muestra que termina. \\
5 & Hilo 2 & El proceso principal continúa, ejecuta \texttt{start()} y el Hilo 2 inicia el procesamiento. \\
6 & Proceso principal & Ejecuta \texttt{hilo2.join()} y espera la terminación del Hilo 2. \\
7 & Hilo 2 & Procesa los datos y muestra que termina. \\
8 & Hilo 3 & El proceso principal ejecuta \texttt{start()}; el Hilo 3 genera el reporte. \\
9 & Proceso principal & \texttt{hilo3.join()} espera la terminación del Hilo 3 y después continúa. \\
\end{tabular}
\end{center}

\emph{¿Qué ocurriría si se eliminaran los \texttt{join()}?} Los tres hilos
podrían iniciar sin un orden garantizado. El Hilo 2 podría intentar usar los
datos antes de que el Hilo 1 los obtenga, y el Hilo 3 podría generar el reporte
antes de que el procesamiento termine. Además, el proceso principal podría
mostrar el mensaje final antes de que las tareas concluyeran. Se perdería la
dependencia entre las tareas.

\espaciofoto{Captura de pantalla 2026-09-16 a la(s) 11.05.01 a.m..png}{Fotografía del código Java y de la salida real del programa ejecutado desde la computadora}{Código Java y salida del programa de sincronización de hilos con \texttt{join()}.}{9cm}

\chapter{Sincronización de procesos con \texttt{waitFor()}}

\textbf{Desarrollen un segundo ejemplo en Java utilizando \texttt{ProcessBuilder}
para iniciar un proceso externo. El proceso principal debe crear e iniciar el
proceso, mostrar que fue iniciado, utilizar \texttt{waitFor()} para esperar su
finalización, mostrar que terminó y continuar con la ejecución del programa.}

\emph{Entradas.} El proceso principal recibe la ruta del ejecutable de Java de
la JVM actual y la ruta del archivo Java que contiene el proceso hijo. También
utiliza un objeto \texttt{ProcessBuilder} para describir el proceso externo.

\emph{Procesamiento.} El proceso principal crea el \texttt{ProcessBuilder},
conecta la entrada y salida del proceso hijo con el terminal e inicia el
proceso mediante \texttt{start()}. Después ejecuta \texttt{waitFor()} y queda
en espera hasta que el proceso hijo termina. Cuando la espera concluye, el
programa principal obtiene el código de salida y continúa con un mensaje final.

\emph{Salidas.} El terminal muestra que el proceso principal crea e inicia el
proceso hijo, que el proceso hijo inicia y termina su tarea, y que el proceso
principal continúa después de \texttt{waitFor()}. El código de salida igual a
cero indica que el proceso hijo terminó correctamente.

\emph{Estructuras de control utilizadas.} Se emplean las clases
\texttt{ProcessBuilder} y \texttt{Process}, los métodos \texttt{start()},
\texttt{inheritIO()} y \texttt{waitFor()}, además de un bloque
\texttt{try-catch} mediante las excepciones declaradas en \texttt{main}.
El proceso hijo utiliza una pausa para hacer visible el periodo en que el
proceso principal permanece esperando.

\emph{Conceptos.} \texttt{ProcessBuilder} permite crear un proceso externo y
\texttt{start()} lo inicia. \texttt{waitFor()} hace que el proceso principal
espere hasta que el proceso hijo termine y devuelve su código de salida. Esta
es una coordinación entre procesos: el proceso principal y el proceso hijo
son ejecuciones separadas, por lo que no comparten directamente las variables
del programa principal.

\emph{Prueba de escritorio.} La siguiente secuencia representa el flujo
solicitado, incluida la espera del proceso principal.

\begin{center}
\scriptsize
\begin{tabular}{c|p{3.2cm}|p{7.8cm}}
Paso & Proceso & Acción y estado \\
\hline
1 & Proceso principal & Crea el objeto \texttt{ProcessBuilder} con el ejecutable y el archivo Java del proceso hijo. \\
2 & Proceso principal & Ejecuta \texttt{start()} y muestra que el proceso hijo fue iniciado. \\
3 & Proceso hijo & Inicia su tarea externa y muestra el mensaje correspondiente. \\
4 & Proceso principal & Ejecuta \texttt{waitFor()} y permanece en espera bloqueada hasta la finalización del proceso hijo. \\
5 & Proceso hijo & Termina su tarea y devuelve el código de salida cero. \\
6 & Proceso principal & \texttt{waitFor()} termina; el proceso principal muestra el código de salida. \\
7 & Proceso principal & Continúa con la ejecución y muestra el mensaje final. \\
\end{tabular}
\end{center}

\emph{¿Qué función cumple \texttt{waitFor()}?} Coordina al proceso principal
con el proceso hijo. El proceso principal no continúa después de esa llamada
hasta que el proceso hijo termina; además, la llamada devuelve el código de
salida del proceso hijo.

\espaciofoto{Captura de pantalla 2026-09-16 a la(s) 11.05.40 a.m..png}{Fotografía del código Java del proceso hijo y de su salida real}{Proceso hijo con \texttt{ProcessBuilder}.}{7cm}

\espaciofoto{Captura de pantalla 2026-09-16 a la(s) 11.06.17 a.m..png}{Fotografía del código Java del proceso principal y de la salida real}{Proceso principal con \texttt{waitFor()}.}{9cm}

\chapter{Análisis y conclusiones}

\emph{¿Qué función cumple \texttt{join()}?} \texttt{join()} establece una
dependencia entre hilos. El hilo que ejecuta la llamada espera la terminación
del hilo indicado y después puede continuar. En el primer programa permite
que la obtención de datos termine antes de iniciar el procesamiento, y que el
procesamiento termine antes de generar el reporte.

\emph{¿Qué diferencia existe entre sincronizar hilos y procesos en los
ejemplos realizados?} Los hilos pertenecen al mismo programa y se ejecutan
dentro de la misma JVM. En el primer ejemplo, \texttt{join()} ordena tareas
realizadas por objetos \texttt{Runnable} y permite que las tareas consecutivas
utilicen los datos compartidos. En el segundo ejemplo, \texttt{ProcessBuilder}
inicia otra ejecución de Java; \texttt{waitFor()} coordina su finalización y
el proceso principal continúa cuando recibe el código de salida.

\emph{Espera.} En ambos casos el hilo o proceso principal debe esperar un
evento: la terminación de otro hilo o la terminación de otro proceso. Esta
espera establece el orden de ejecución requerido y evita que una tarea
dependiente continúe antes de tiempo.

\emph{Conclusión.} Los dos programas muestran que la sincronización coordina
la ejecución de procesos e hilos para que las operaciones ocurran en el orden
adecuado. \texttt{start()} permite iniciar una nueva ejecución, mientras que
\texttt{join()} y \texttt{waitFor()} establecen el punto en el que la
ejecución debe esperar. En el primer programa, eliminar \texttt{join()} haría
que el orden de las tareas dependiera del planificador y podría producir un
reporte sin datos procesados. En el segundo programa, eliminar \texttt{waitFor()}
permitiría que el proceso principal continuara antes de que terminara el
proceso hijo.

\cleardoublepage
\phantomsection
\addcontentsline{toc}{chapter}{Referencias}
\nocite{*}
\bibliography{template/references}

\end{document}
